\section{Permanence properties}
\label{I.9}

Let $A\to B$ be a local \'etale homomorphism. We examine here
some cases in which a certain property of $A$ implies the
same property for~$B$, or conversely.
% original p. 16
A number of such propositions already follow merely from
the fact that $B$ is \emph{quasi-finite} and \emph{flat}
over~$A$, and we shall confine ourselves to ``recalling''
a few of them: $A$ \emph{and} $B$ \emph{have the same Krull
dimension and the same depth} (Serre's ``cohomological codimension'',
in the terminology still current). It follows, for example,
that \emph{$A$ is Cohen--Macaulay if and only if $B$ is so}.
Moreover, for every prime ideal $\goth{q}$ of~$B$, inducing
$\goth{p}$ on~$A$, $B_{\goth{q}}$ will again be quasi-finite
and flat over~$A_{\goth{p}}$, provided that $B$ is assumed
to be a localization of an algebra of finite type over~$A$
(this follows from the fact that the set of points where
a morphism of finite type is quasi-finite \resp flat is
open); furthermore, \emph{every} prime ideal $\goth{p}$ of~$A$
is induced by a prime ideal $\goth{q}$ of~$B$ (since $B$ is
\emph{faithfully} flat over~$A$). It follows, for example,
that \emph{$\goth{p}$ and~$\goth{q}$ have the same height};
and also that \emph{$A$ has no embedded prime ideals if
and only if~$B$ has none}.

We shall therefore restrict ourselves to propositions more
specific to the case of \'etale morphisms.

\begin{proposition}
\label{I.9.1}
Let $A\to B$ be a local \'etale homomorphism. For $A$ to be
regular, it is necessary and sufficient that $B$ be so.
\end{proposition}

Indeed, let $k$ be the residue field of~$A$, and $L$ that
of~$B$. Since $B$ is flat over~$A$ and $L=B\otimes_A k$,
\ie $\goth{n}=\goth{m}B$ (where $\goth{m}$,~$\goth{n}$
are the maximal ideals of~$A$,~$B$), the $\goth{m}$-adic
filtration on~$B$ is identical to its $\goth{n}$-adic
filtration, and one has
\[
\gr^*(B)=\gr^*(A)\otimes_k L.
\]
It follows that $\gr^*(B)$ is a polynomial algebra over~$L$
if and only if $\gr^*(A)$ is a polynomial algebra over~$k$,
qed. (N.B.\ The fact that $L/k$ is separable has not been used.)

\begin{corollary}
\label{I.9.2}
Let $f\colon X\to Y$ be an \'etale morphism. If $Y$ is
regular, so is $X$; the converse holds if $f$ is surjective.
\end{corollary}

% The corrected source repeats 9.2. Preserve it with a distinct PDF anchor.
\setcounter{theorem}{1}
\begingroup
\renewcommand{\theHtheorem}{\theHsection.\arabic{theorem}.bis}
\begin{proposition}
\label{prop:I.9.2}
Let $f\colon X\to Y$ be an \'etale morphism. If $Y$ is
reduced, so is~$X$; the converse holds if $f$ is surjective.
\end{proposition}
\endgroup

This is equivalent to the following

\begin{corollary}
\label{I.9.3}
Let $f\colon A\to B$ be a local \'etale homomorphism, with
$B$ isomorphic to an
% original p. 17
$A$-algebra which is a localization of an $A$-algebra of
finite type. For $A$ to be reduced, it is necessary and
sufficient that $B$ be so.
\end{corollary}

Necessity is trivial, since $A\to B$ is injective ($B$ being
faithfully flat over~$A$). Sufficiency: let $\goth{p}_i$ be
the minimal prime ideals of~$A$; by hypothesis the natural
map $A\to\prod_i A/\goth{p}_i$ is injective, so tensoring
with the flat $A$-module $B$, one finds that
$B\to\prod_i B/\goth{p}_iB$ is injective, and one is reduced
to proving that the $B/\goth{p}_iB$ are reduced. Since
$B/\goth{p}_iB$ is \'etale over~$A/\goth{p}_i$, one is
reduced to the case where $A$ is an integral domain. Let~$K$
be its field of fractions; since $A\to K$ is injective,
so is $B\to B\otimes_A K$ ($B$ being $A$-flat), and one is
reduced to proving that the latter ring is reduced. Now,
$B$, being a localization of an $A$-algebra of finite type
over~$A$, is the local ring of a point~$x$ of a scheme
$X=\Spec(C)$ of finite type and \emph{\'etale} over
$Y=\Spec(A)$; hence $B\otimes_A K$ is a localization (with
respect to a suitable multiplicatively closed set) of the
ring $C\otimes_A K$ of~$X\otimes_A K$. Since $X\otimes_A K$
is \'etale over~$K$, its ring is a finite product of fields
(separable extensions of~$K$); hence the same holds
for~$B\otimes_A K$, qed.

\begin{corollary}
\label{I.9.4}
Let $f\colon A\to B$ be a local \'etale homomorphism, and
suppose that $A$ is analytically reduced (\ie the completion
$\hat{A}$ of~$A$ has no nilpotent elements). Then $B$ is
analytically reduced, and a fortiori reduced.
\end{corollary}

Indeed, $\hat{B}$ is \emph{finite} and \'etale over~$\hat{A}$,
and one applies~\ref{I.9.3}.

\begin{theorem}
\label{I.9.5}
Let $f\colon A\to B$ be a local homomorphism, with $B$
isomorphic to a localization of an $A$-algebra of finite
type. Then:
\begin{enumerate}
\item[(i)] If $f$ is \'etale, $A$ is normal if and only if
$B$ is so.
\item[(ii)] If $A$ is normal, $f$ is \'etale if and only if
$f$ is injective and net (and then $B$ is normal by \textup{(i)}).
\end{enumerate}
\end{theorem}

We shall give two different proofs of (i). The first uses
some of the properties of flat quasi-finite morphisms
(recalled at the beginning of this no.) without using
\ref{I.7.6} (and thereby the Main Theorem); the reverse
is true of the second proof. Finally, for~(ii) it seems
that one needs the Main Theorem in any case.

\subsubsection*{First proof}
One uses the following necessary and sufficient condition
% original p. 18
for normality of a noetherian local ring $A$ of dimension
$\neq0$.

\begin{namedstatement}{Serre's criterion}
\textup{(i)}~For every prime ideal $\goth{p}$ of~$A$ of
height~$1$, $A_{\goth{p}}$ is normal (or, equivalently,
regular); \textup{(ii)}~For every prime ideal $\goth{p}$
of~$A$ of height~$\geq2$, the depth of $A_{\goth{p}}$ is
$\geq2$.\footnote{\Cf EGA~IV~5.8.6.}
\end{namedstatement}

We shall admit this criterion here; it is supposed to appear
in the paragraph on flat morphisms. Its main advantage
is that it does not assume a priori that $A$ is reduced,
let alone an integral domain. Here one may already suppose
that $\dim{A}=\dim{B}\neq0$.

By the facts recalled at the beginning of this no., the
prime ideals $\goth{p}$ of~$A$ of height~$1$ (\resp of
height~$\ge2$) are precisely the contractions to~$A$ of
the prime ideals~$\goth{q}$ of~$B$ of height~$1$ (\resp
of height~$\ge2$). Finally, if $\goth{p}$ and $\goth{q}$
correspond, $B_{\goth{q}}$ is \'etale over~$A_{\goth{p}}$,
hence has the same depth as $A_{\goth{p}}$, and is regular
if and only if $A_{\goth{p}}$ is so (\ref{I.9.1}). Applying
Serre's criterion, one finds that $A$ is normal if and only
if $B$ is so.

\subsubsection*{Second proof}
Suppose that $B$ is normal; let $L$ be its field of fractions,
and $K$ that of $A$ ($A$ is an integral domain since $B$ is
so). We saw in the proof of~\ref{I.9.3} that $B\otimes_A K$
is a finite product of fields; since it is contained in $L$
it is a field, and since it contains $B$ it is $L$. An
element of~$K$ integral over~$A$ is integral over~$B$, hence
belongs to $B$ since $B$ is normal, and therefore to $A$
since $B\cap K=A$ (as follows from the fact that $B$ is
faithfully flat over~$A$).

Suppose now that $A$ is normal; let us prove that $B$ is
so. By~\ref{I.7.6}, one has $B=B^\prime_{\goth{n}}$, where
$B^\prime=A[t]/FA[t]$, with $F$ and $\goth{n}$ as
in~\ref{I.7.6}. Thus $L=B\otimes_A K$ will be a localization
of~$B^\prime\otimes_A K=K[t]/FK[t]$, and a product of fields,
finite separable extensions of~$K$; this product (as whenever
one localizes an artinian ring, here $B^\prime_K$, with
respect to a multiplicatively closed set) is a direct
factor of~$B^\prime_K$, thus corresponding to a decomposition
$F=F_1F_2$ in $K[t]$, with the generator of~$L$ corresponding
to $t$ already annihilated by~$F_1$. Now, \emph{since $A$
is normal, the $F_i$ belong to $A[t]$} (assuming them monic).
Observing that $B\to L=B\otimes_A K$ is injective (since
$A\to K$ is so and $B$ is flat over~$A$), it follows that
one already has $F_1(u)=0$, with $u$ the class of~$t$ in~$L$.
Assuming that $F$ was chosen of minimal degree, it follows
that $F_2=1$ (N.B.\ One has
$F^\prime(u)=F^\prime_1(u)F_2(u)+F_1(u)F^\prime_2(u)
=F^\prime_1(u)F_2(u)$ since $F_1(u)=0$, whence
$F^\prime_1(u)\neq0$ since $F^\prime(u)\neq0$.)

Thus
% original p. 19
one has
\begin{equation*}
\label{eq:I.9.5.*}
\tag{$*$} L=B\otimes_A K=K[t]/FK[t]
\end{equation*}
and $F$ is consequently a separable polynomial in $K[t]$
(but of course not necessarily in $A[t]$). (N.B.\ For the
moment, we have essentially shown only that in~\ref{I.7.6}
one can choose $F$ and $\goth{n}$ in such a way that---with
the notation used here---$B^\prime\to B^\prime_{\goth{n}}=B$
is \emph{injective}; we used the normality of $A$ for this;
I do not know whether it remains true without the normality
assumption.)

Let us now recall the well-known lemma taken from Serre's
course of last year:

\begin{lemma}
\label{I.9.6}
Let $K$ be a ring, $F\in K[t]$ a monic separable polynomial,
$L=K[t]/FK[t]$, and $u$ the class of~$t$ in~$L$ (so that
$F^\prime(u)$ is an invertible element of~$L$). Then one
has the formulas (where $n=\deg{F}$):
\begin{enumerate}
\item[] $\trace_{L/K}u^i/F^\prime(u)=0$\quad if\quad $0\le i<n-1$,
\item[] $\trace_{L/K}u^{n-1}/F^\prime(u)=1$.
\end{enumerate}
\end{lemma}

\begin{corollary}
\label{I.9.7}
The determinant of the matrix
$(u^j\cdot u^i/F^\prime(u))_{0\le i,j\le n-1}$ is equal
to $(-1)^{n(n-1)/2}$, hence is invertible in every subring
$A$ of~$K$.
\end{corollary}

\begin{corollary}
\label{I.9.8}
Let $A$ be a subring of~$K$, $V$ the $A$-module generated
by the $u^i$ ($0\le i\le n-1$) in $L$, and $V^\prime$
the $A$-submodule of~$L$ consisting of the $x\in L$ such
that $\trace_{L/K}(xy)\in A$ for every $y\in V$ (\ie for
$y$ of the form $u^i$, $0\le i\le n-1$). Then $V^\prime$
is the $A$-module with basis the $u^i/F^\prime(u)$
($0\le i\le n-1$).
\end{corollary}

\begin{corollary}
\label{I.9.9}
Suppose that $K$ is the field of fractions of a normal
integral domain $A$, with $F$ having its coefficients
in $A$. Then, with the notation of~\ref{I.9.8}, $V^\prime$
contains the integral closure $A^\prime$ of~$A$ in $L$,
which is therefore contained in $A[u]/F^\prime(u)$ and
a fortiori in $A[u][F^\prime(u)^{-1}]$.
\end{corollary}

Let us apply this last corollary to the situation obtained
in the proof: since $F^\prime(u)$ is invertible in $B$,
which contains $A[u]$, $B$ contains $A^\prime$. By the
Main Theorem (or from the fact that $B=A[u]_{\goth{n}}$),
$B$ is a localization of $A^\prime$. Since $A^\prime$ is
normal, so is~$B$.

\subsubsection*{Proof of \textup{(ii)}}
One
% original p. 20
proceeds as in the preceding proof to show that, in
\ref{I.7.6}, one can choose $F$ so that ($*$) still holds.
The only obstacle a priori is that, since $B$ is no longer
assumed flat over~$A$, one can no longer assert that
$B\to L$ is injective, so that the argument applies a
priori only to the image $B_1$ of~$B$ under the said
homomorphism. It follows at once that $B_1$ is flat
over~$A$ (being a localization of a free algebra over~$A$).
By~\ref{I.4.8}, the morphism $B\to B_1$ is \'etale,
hence an isomorphism, which completes the proof.

(From the point of view of exposition, one should
interchange the last two proofs and put the formal
calculations of the lemma and its corollaries in a
separate no.)

\begin{corollary}
\label{I.9.10}
Let $f\colon X\to Y$ be an \'etale morphism. If $Y$ is
normal, so is $X$; the converse holds if $f$ is surjective.
\end{corollary}

\begin{corollary}
\label{I.9.11}
Let $f\colon X\to Y$ be a dominant morphism, with $Y$
normal and $X$ connected. If $f$ is net, $f$ is \'etale,
hence $X$ is normal and consequently (being connected)
irreducible.
\end{corollary}

Let $U$ be the set of points where $f$ is \'etale; it is
open, and it suffices to show that it is also closed and
nonempty. $U$ contains the inverse image of the generic
point of~$Y$ (since, for an algebra over a field,
unramified $=$ \'etale), hence ($X$ dominating $Y$) is
nonempty. If $x$ belongs to the closure of~$U$, then it
belongs to the closure of an irreducible component $U_i$
of~$U$, hence to an irreducible component
$X_i\overset{v}{=}\bar{U}_i$ of~$X$ which meets $U$ and
consequently dominates $Y$ (since every component of~$U$,
flat over~$Y$, dominates $Y$). Thus, if $y$ is the
projection of~$x$ on~$Y$, $\cal{O}_y\to\cal{O}_x$ is
\emph{injective} (taking into account that $\cal{O}_y$
is an integral domain). Since $\cal{O}_y$ is normal
and $\cal{O}_y\to\cal{O}_x$ is net, one concludes
using~\ref{I.9.5}(ii).

\begin{corollary}
\label{I.9.12}
Let $f\colon X\to Y$ be a dominant morphism of finite
type, with $Y$ normal and $X$ irreducible. Then the set
of points where $f$ is \'etale is identical to the
complement of the support of $\Omega^1_{X/Y}$, \ie to
the complement of the subprescheme of~$X$ defined by
the different ideal $\goth{d}_{X/Y}$.
\end{corollary}

(This
% original p. 21
is the ``less trivial'' statement alluded to in the
remark in~\No\ref{I.4}.)

\begin{remark}
One should beware of believing that a connected \'etale
covering of an irreducible scheme is itself irreducible
when the base is not assumed normal. This question
will be studied in~\No\ref{I.11}.
\end{remark}
